% ============================================================
%  Hiero Workflow App — Architecture-First Response
%  2-PAGE EDITION  |  Darshit  |  LFDT Mentorship 2026
% ============================================================
\documentclass[9pt, a4paper]{article}

% ── Encoding ────────────────────────────────────────────────
\usepackage[T1]{fontenc}
\usepackage[utf8]{inputenc}

% ── Tight geometry ──────────────────────────────────────────
\usepackage[
  top        = 1.0cm,
  bottom     = 0.85cm,
  left       = 0.9cm,
  right      = 0.9cm,
  headheight = 9pt,
  headsep    = 4pt,
  footskip   = 7pt,
]{geometry}

% ── Spacing ─────────────────────────────────────────────────
\usepackage{parskip}
\setlength{\parindent}{0pt}
\setlength{\parskip}{1pt}
\raggedbottom

% ── Two-column layout ───────────────────────────────────────
\usepackage{multicol}
\setlength{\columnsep}{0.42cm}
\setlength{\columnseprule}{0.3pt}
\def\columnseprulecolor{\color{hieroRule}}
\setlength{\multicolsep}{3pt}

% ── Tables ──────────────────────────────────────────────────
\usepackage{booktabs}
\usepackage{tabularx}
\usepackage{array}
\newcolumntype{L}[1]{>{\raggedright\arraybackslash}p{#1}}
\renewcommand{\arraystretch}{1.0}

% ── Colours ─────────────────────────────────────────────────
\usepackage{xcolor}
\definecolor{hieroBlue}{HTML}{1A3A5C}
\definecolor{hieroAccent}{HTML}{2E75B6}
\definecolor{hieroGray}{HTML}{5A5A5A}
\definecolor{hieroRule}{HTML}{BBBBBB}
\definecolor{warnAmber}{HTML}{FFF3CD}
\definecolor{warnBorder}{HTML}{B8860B}
\definecolor{invariantBg}{HTML}{E8F4E8}
\definecolor{invariantBorder}{HTML}{2E7D32}
\definecolor{thesisBg}{HTML}{EAF2FB}
\definecolor{thesisBorder}{HTML}{1A3A5C}
\definecolor{tagK}{HTML}{1565C0}
\definecolor{tagR}{HTML}{2E7D32}
\definecolor{tagA}{HTML}{6A1B9A}
\definecolor{tagE}{HTML}{B71C1C}

% ── Hyperlinks ──────────────────────────────────────────────
\usepackage{hyperref}
\hypersetup{
  colorlinks = true,
  linkcolor  = hieroAccent,
  urlcolor   = hieroAccent,
  pdftitle   = {Hiero Workflow App — Architecture-First Response},
  pdfauthor  = {Darshit | LFDT Mentorship 2026},
}

% ── Framed boxes (ultra-tight) ──────────────────────────────
\usepackage{mdframed}

\newmdenv[backgroundcolor=thesisBg, linecolor=thesisBorder, linewidth=1pt,
  leftline=true, rightline=false, topline=false, bottomline=false,
  innerleftmargin=5pt, innerrightmargin=4pt,
  innertopmargin=2pt, innerbottommargin=2pt,
  skipabove=2pt, skipbelow=1pt]{thesisbox}

\newmdenv[backgroundcolor=invariantBg, linecolor=invariantBorder, linewidth=0.6pt,
  innerleftmargin=4pt, innerrightmargin=4pt,
  innertopmargin=1pt, innerbottommargin=2pt,
  skipabove=2pt, skipbelow=1pt]{invariantbox}

\newmdenv[backgroundcolor=warnAmber, linecolor=warnBorder, linewidth=0.6pt,
  innerleftmargin=4pt, innerrightmargin=4pt,
  innertopmargin=2pt, innerbottommargin=2pt,
  skipabove=2pt, skipbelow=1pt]{warnbox}

% ── Inline criterion tag ────────────────────────────────────
%  Renders as a tiny coloured badge inline — zero extra vertical space
\newcommand{\ic}[2]{%
  {\colorbox{tag#1}{\color{white}\bfseries\fontsize{5.5}{6.5}\selectfont\,[#1\,=\,#2]\,}}%
}

% ── Marks & link helpers (safe in section titles) ───────────
\DeclareRobustCommand{\mb}[1]{%
  \texorpdfstring{\hfill{\fontsize{6.5}{7}\selectfont\color{hieroGray}[#1]}}{}%
}
\DeclareRobustCommand{\slink}[1]{%
  \texorpdfstring{\hspace{4pt}{\fontsize{7}{8}\selectfont\color{hieroAccent}\textbf{#1}}}{}%
}

% ── Lists (ultra-tight) ─────────────────────────────────────
\usepackage{enumitem}
\setlist[itemize]  {leftmargin=1.1em, itemsep=0pt, topsep=1pt, parsep=0pt}
\setlist[enumerate]{leftmargin=1.1em, itemsep=0pt, topsep=1pt, parsep=0pt}

% ── Section headings ────────────────────────────────────────
\usepackage{titlesec}
\titleformat{\section}
  {\small\bfseries\color{hieroBlue}}{Q\thesection.}{0.4em}{}
  [\vspace{-3pt}\color{hieroRule}\rule{\linewidth}{0.3pt}]
\titlespacing*{\section}{0pt}{4pt}{1pt}

% ── Header / footer ─────────────────────────────────────────
\usepackage{fancyhdr}
\pagestyle{fancy}
\fancyhf{}
\renewcommand{\headrulewidth}{0.3pt}
\renewcommand{\footrulewidth}{0pt}
\fancyhead[L]{\fontsize{7}{8}\selectfont\color{hieroGray}%
  \textit{Hiero Workflow App — Architecture-First Response}}
\fancyhead[R]{\fontsize{7}{8}\selectfont\color{hieroGray}LFDT Mentorship 2026}
\fancyfoot[R]{\fontsize{7}{8}\selectfont\color{hieroGray}\thepage}

% ── Table font shortcut ─────────────────────────────────────
\newcommand{\tsize}{\fontsize{7.8}{9.5}\selectfont}

% ============================================================
\begin{document}
% ============================================================
\thispagestyle{fancy}

% ── Cover ───────────────────────────────────────────────────
\begin{center}
  {\normalsize\bfseries\color{hieroBlue}Hiero Workflow App — Architecture-First Response}\\[1pt]
  {\fontsize{7.5}{9}\selectfont\color{hieroGray}%
    Pre-interview task \quad|\quad Darshit \quad|\quad LFDT Mentorship 2026}
\end{center}

\vspace{1pt}{\color{hieroRule}\hrule}\vspace{2pt}

\begin{thesisbox}
\fontsize{8}{9.5}\selectfont
\textbf{Thesis.}\ Build a shared automation core in \texttt{sdk-automations}, expose it via a
\textbf{GitHub Action adapter first}, retain Probot\,/\,GitHub App as the later orchestration
path. SDKs keep local workflow ownership, security controls, and policy; reusable decision
logic moves into the shared core.\\[1pt]
\textbf{Centralise:} decision logic, schemas, fixtures, parity tests, adapters, releases.\quad
\textbf{Keep local:} workflow YAML, triggers, permissions, checkout, harden-runner, policy.\quad
\textbf{Never rewrite:} SDK security controls or maintainer governance.
\end{thesisbox}

\vspace{1pt}{\color{hieroRule}\hrule}\vspace{2pt}

% ── Body — two columns ──────────────────────────────────────
\begin{multicols}{2}
\fontsize{8.5}{10.5}\selectfont

% ============================================================
\section{Well-Architected Principles Applied to Hiero \mb{4 pts — K=2, R=2}}
% ============================================================

\ic{K}{2}\ \textbf{Clear responsibility.}
SDK repos own triggers, permissions, checkout, harden-runner, secrets, and local policy.
The central repo owns automation decisions, schemas, adapters, fixtures, and releases.
No step hides a security control or merges two governance domains.

\ic{K}{2}\ \textbf{Secure by design.}
Least-privilege \texttt{GITHUB\_TOKEN}; no untrusted PR code in privileged jobs;
pin actions by full commit SHA; CODEOWNERS on config; fail closed on invalid config.

\ic{R}{2}\ \textbf{Low coupling\,/\,high cohesion.}
Automation rules live in \texttt{packages/core}; GitHub Actions and Probot are thin adapters
over the same behaviour — no duplication of logic.

\ic{R}{2}\ \textbf{Observable and testable.}
Every automation supports dry-run output, structured logs, idempotent writes, golden
fixtures, and parity tests before write mode is enabled upstream.

\textbf{Evolvable.}
New automations are registered in the core and versioned via config schema; no
cross-SDK script copying required.

% ============================================================
\section{Proposed Final System Architecture \mb{4 pts — R=2, A=2} \slink{↗ Link}}
% ============================================================

\ic{R}{2}\ic{A}{2}\ Four zones: \textbf{Python\,/\,C++ SDK repos} (own all workflow control and
security) → \textbf{GitHub Action adapter} (current path; validates inputs, calls core,
supports dry-run\,/\,writes) → \textbf{\texttt{packages/core}} (registry, review-sync, assignment,
PR checks, comment builders, label decisions, fixtures, parity tests) → \textbf{GitHub API}.
Probot\,/\,GitHub App calls the same core later — no second behaviour implementation.

\begin{invariantbox}
\fontsize{7.5}{9}\selectfont
\textbf{Invariant:} centralise reusable decision logic; do not centralise repository control
until governance, security, and parity gates are satisfied.
\end{invariantbox}

% ============================================================
\section{Architectural Boundaries and Why They Exist \mb{4 pts — R=2, E=2}}
% ============================================================

\ic{R}{2}

\vspace{1pt}
{\tsize
\begin{tabularx}{\linewidth}{@{}L{1.85cm}X@{}}
\toprule
\textbf{Boundary} & \textbf{What lives there and why} \\
\midrule
Repo-local &
  YAML, permissions, checkout, harden-runner, concurrency, token. Remain local because these
  are security and maintainer-ownership controls. Moving them hides the security surface
  from the team accountable for it. \\
Shared-core &
  Label decisions, assignment, review-sync, PR checks, comment builders, retry helpers,
  schemas, fixtures. Central because duplication creates maintenance cost and behavioural drift. \\
Policy &
  Labels, team handles, docs links, limits, enabled automations. Protected local config so
  SDK-specific choices do not become code forks in the shared repo. \\
Hosted-app &
  Probot hosting deferred: introduces webhook secrets, installation permissions, availability
  SLAs, incident response, and operational ownership not yet warranted. \\
\bottomrule
\end{tabularx}}

\vspace{1pt}
\begin{invariantbox}
\fontsize{7.5}{9}\selectfont
\ic{E}{2}\ \textbf{Key evaluation:} Each boundary maps to a real risk if crossed — hiding
permissions, creating hidden code forks, or adding unowned operational surface.
\end{invariantbox}

% ─── column break: Q4 starts in right column ───────────────
\columnbreak

% ============================================================
\section{Tooling Choices \mb{6 pts — R=2, A=2, E=2}}
% ============================================================

\ic{R}{2}

\vspace{1pt}
{\tsize
\begin{tabularx}{\linewidth}{@{}L{1.65cm}X@{}}
\toprule
\textbf{Area} & \textbf{Recommendation and justification} \\
\midrule
Caller runtime &
  \textbf{GitHub Actions first.} Matches existing SDK ops; maintainers keep YAML, permissions,
  checkout, concurrency, and harden-runner local. No new hosting in Phase 1. \\
Core language &
  \textbf{Node.js\,/\,Octokit.} Native to GitHub automation; bundles cleanly as an action
  artefact; no secondary runtime for the adapter layer. \\
Shared logic &
  \textbf{\texttt{packages/core}.} Adapter-agnostic modules. Actions and Probot adapters call
  the same functions — behaviour tested once, delivered twice. \\
Config &
  \textbf{YAML\,/\,JSON + JSON Schema.} Versioned contract; invalid config fails closed before
  any write occurs. \\
Verification &
  \textbf{Unit + mock + fixtures + canaries + parity.} Each layer catches a different failure
  class: logic, API contracts, or real-world call paths. \\
Security ops &
  \textbf{Harden-Runner, CODEOWNERS, SHA pins, Dependabot, rollback docs.}
  Governance gate, not optional extras. \\
\bottomrule
\end{tabularx}}

\begin{invariantbox}
\fontsize{7.5}{9}\selectfont
\ic{A}{2}\ic{E}{2}\ Every choice is justified by what it \emph{prevents}. Node.js bundles
natively; JSON Schema enforces a contract. Tooling that cannot be pinned, audited, or
rolled back is not acceptable in a security-sensitive automation path.
\end{invariantbox}

% ============================================================
\section{How This Solution Builds Iteratively \mb{2 pts — A=2}}
% ============================================================

\ic{A}{2}
\begin{enumerate}
  \item \textbf{Architecture + reversible work in parallel.} Core skeleton, schema, fixtures,
        and dry-run path proceed now. Production writes wait for architecture gates.
  \item \textbf{Pilot Python review-sync first.} Bounded, canary-proven. One controlled write
        mode enabled after parity tests pass.
  \item \textbf{Map assignment logic before extracting it.} Same filenames do not guarantee
        identical policy; extraction without a behaviour map creates hidden drift.
  \item \textbf{C++ follows investigation.} Understand awkward workflows before introducing
        an unproven wrapper.
  \item \textbf{Probot comes after core is trusted} — as an adapter over the validated
        contract, not a second implementation.
\end{enumerate}

% ============================================================
\section{Malicious Pull Request Risks and Safeguards \mb{10 pts — K=4, R=3, A=3}}
% ============================================================

\ic{K}{4}\ Highest-risk design area. The architecture assumes a malicious contributor controls:
PR title\,/\,body, branch name, changed files, artefacts from untrusted code, and
workflow\,/\,config changes in the PR branch.

\vspace{1pt}
{\tsize
\begin{tabularx}{\linewidth}{@{}L{1.9cm}X@{}}
\toprule
\textbf{Risk} & \textbf{Safeguard} \\
\midrule
\texttt{pull\_request\_target} abuse &
  Never checkout PR head in privileged jobs; load config from default branch; use
  job-level least privilege; avoid secrets in PR contexts. \\
Token\,/\,secret exfiltration &
  Harden-runner early in every job; restrict\,/\,audit egress; mask sensitive values;
  review logs during all canary runs. \\
Action\,/\,dependency compromise &
  Pin all actions by full commit SHA; bundle central actions; use Dependabot and
  \texttt{npm audit} continuously. \\
Config\,/\,policy tampering &
  CODEOWNERS on \texttt{.github/hiero-automation.*}; schema-validate every invocation;
  reject unknown fields; fail closed on any error. \\
Artefact poisoning &
  Validate artefact identity: run source, run ID, path, checksum. Never treat artefacts
  from untrusted PR code as an authority signal. \\
Bot loops\,/\,broad writes &
  Actor\,/\,bot guards, idempotent markers, scoped label allowlists, concurrency keys,
  rollback by disabling the caller workflow or pinning back a SHA. \\
\bottomrule
\end{tabularx}}

\vspace{1pt}
\begin{warnbox}
\fontsize{7.5}{9}\selectfont
\ic{R}{3}\ic{A}{3}\ \textbf{Systemic safeguards:} validate actor type, event type, repo,
PR number, label allowlists, and config schema \emph{before} acting — unknown config
fails closed. All writes are idempotent and scoped to known bot markers only.
For \texttt{workflow\_run}: validate artefact identity, run source, and expected paths
locally before the action consumes anything. For Probot (later): verify webhook signatures,
use installation-scoped permissions, store secrets outside repos, implement rate-limit
handling, and provide a clear disable\,/\,rollback path.
\end{warnbox}

% ============================================================
\section{Configuration Standardisation vs.\ User Input \mb{6 pts — K=2, R=2, A=2}}
% ============================================================

\ic{K}{2}\ic{R}{2}

\vspace{1pt}
{\tsize
\begin{tabularx}{\linewidth}{@{}XX@{}}
\toprule
\textbf{Standardise centrally} & \textbf{Allow per-repo (with guardrails)} \\
\midrule
Schema version, automation names, dry-run\,/\,write contract, result\,/\,log format,
bot comment markers, retry logic, error categories, release compatibility, permission docs.
&
Label names, team handles, docs links, support channels, assignment limits, enabled
automations, schedules, branch filters, repo-specific message links. \\
\bottomrule
\end{tabularx}}

\begin{invariantbox}
\fontsize{7.5}{9}\selectfont
\ic{A}{2}\ \textbf{Rule of thumb:} common behaviour $\to$ core; repo policy $\to$ protected
config validated by the central schema; one-off workflow\,/\,security constraint $\to$ keep
strictly local. Avoid arbitrary scripts, unlimited switches, and template overrides —
those recreate the maintenance tax centrally.
\end{invariantbox}

% ============================================================
\section{Community Leverage Plan \mb{10 pts — K=3, R=3, A=4}}
% ============================================================

\ic{K}{3}\ \textbf{Public issue backlog.} Architecture becomes a contribution surface
covering behaviour mapping, fixture collection, docs, config examples, golden tests, canary
workflows, security review, and SDK adoption guides — useful work that does not require
maintainers to triage risky rewrites.

\ic{R}{3}\ \textbf{Skill-based labels:} \texttt{good-first-issue} (docs\,/\,fixtures)
$\cdot$ \texttt{beginner} (tests\,/\,samples) $\cdot$ \texttt{intermediate}
(adapters\,/\,config) $\cdot$ \texttt{advanced} (security pilots\,/\,C++ investigation).

\ic{A}{4}\ \textbf{RFC per automation migration.} Each RFC states: reusable logic,
repo-local boundary, policy config, fixtures, rollback plan, and owner. No production
behaviour change without dry-run proof, parity tests, a security checklist, and maintainer
approval. Small, reviewable PRs only.

\vspace{1pt}
{\tsize
\begin{tabularx}{\linewidth}{@{}L{1.85cm}X@{}}
\toprule
\textbf{Work package} & \textbf{Good community issue shape} \\
\midrule
Behaviour mapping &
  Compare Python and C++ for one automation; record common logic, policy differences,
  fixtures, and unknowns. \\
Fixture\,/\,parity tests &
  Convert real examples to golden cases so central behaviour is proved before writes. \\
Docs\,/\,examples &
  Caller workflow examples, config samples, rollback instructions, security checklists. \\
Canary support &
  Run dry-run pilots, collect logs, identify false positives, file follow-up issues. \\
Security review &
  Threat-model \texttt{pull\_request\_target}, artefact flows, permissions, installation
  scopes, and release pinning. \\
\bottomrule
\end{tabularx}}

\medskip
\textbf{Demos and office hours.} Publish canary reproduction steps so contributors can
verify call paths independently. Python\,/\,C++ maintainers review for behaviour; central
maintainers review for schema, release, and security.

% ============================================================
\section{Architecture vs.\ Immediate Development \mb{4 pts — R=2, E=2}}
% ============================================================

\ic{R}{2}

\vspace{1pt}
{\tsize
\begin{tabularx}{\linewidth}{@{}L{1.55cm}L{0.85cm}X@{}}
\toprule
\textbf{Option} & \textbf{Risk} & \textbf{Assessment} \\
\midrule
Skip architecture &
  High &
  Replicates existing problem: independent evolution, inconsistent security, no parity baseline. \\
Full freeze &
  Medium &
  Delays reversible work unnecessarily; fixtures, schema, and dry-run adapters carry no risk. \\
\textbf{2-week arch + reversible work} &
  \textbf{Low} &
  \textbf{Recommended.} Architecture gates production writes; reversible work runs in parallel. \\
\bottomrule
\end{tabularx}}

\begin{invariantbox}
\fontsize{7.5}{9}\selectfont
\ic{E}{2}\ \textbf{Allocation:}
\textbf{Wk\,1} — context diagram, boundary definitions, threat model, config schema contract.
\textbf{Wk\,2} — behaviour maps, pilot criteria, rollback plan, community issue backlog.\\
\textbf{Safe to begin now:} fixtures, schema validation, core modules, dry-run adapters, docs.
\textbf{Must wait:} upstream write mode, SDK security control changes, C++ wrapper introduction.
\end{invariantbox}

\end{multicols}

\vspace{2pt}
{\color{hieroRule}\hrule}
\vspace{2pt}
{\fontsize{6.8}{8}\selectfont\color{hieroGray}\centering
  References: arc42.org $\cdot$ GitHub Actions secure-use docs $\cdot$ StepSecurity Harden-Runner
  $\cdot$ LFDT Issue \#73 $\cdot$ C++ Issue \#1627 $\cdot$ Richards, \textit{Software Architecture Patterns}
  (O'Reilly, 2015)\par}

% ============================================================
\end{document}
% ============================================================